\documentclass[12pt]{article}
\usepackage[utf8]{inputenc}
\usepackage[T1]{fontenc}
\usepackage{hyperref}
\title{Tombonuo Language Resources for K-12 Curriculum Development}
\author{Compiled by Claude Sonnet 5 (claude-sonnet-5) for Andr\'as Kornai}
\date{}
\begin{document}
\maketitle

\section{Introduction}
Tombonuo (Tombonuwo, Tambonuo; also Sungai, Lobu, Mutangar Tombonuo; ISO 639-3: txa) is a Paitanic language of northwest Sabah, Malaysia (Pitas and Labuk-Sugut districts). Speaker figures vary (about 10,000 versus 3,000 in one cited count). Sources disagree on whether it is written: Omniglot describes a Latin alphabet, while Ethnologue (via the Endangered Languages entry) lists it as unwritten.

\section{Primary Educational Resources}
No primers, textbooks or school programmes were found in search results.

\section{Linguistic and Reference Materials}
\begin{itemize}
\item Tombonuwo language overview: \url{https://en.wikipedia.org/wiki/Tombonuwo_language}
\item Omniglot page with alphabet and sample text: \url{https://www.omniglot.com/writing/tombonuwo.htm}
\item Numeral survey (Lingkabau dialect): \url{https://lingweb.eva.mpg.de/channumerals/Tombonuo.htm}
\item Library of Congress subject record, which cites a King wordlist (1990): \url{https://id.loc.gov/authorities/subjects/sh93006465.html}
\item Endangered Languages / Ethnologue entry: \url{https://endangeredlanguages.com/node/89736}
\end{itemize}
Julie King's 1984 \emph{The Paitanic language family} is cited by these sources, but I did not find a link to it.

\section{Religious and Community Resources}
Tombonuo New Testament (2002, \emph{Tinongaran Nu Kinoringan}, Bible Society of Malaysia), with audio, at Digital Bible Society: \url{https://dbs.org/bibles/TXATBL}

\section{Digital and Audio Resources}
The Digital Bible Society page above offers the audio New Testament. Omniglot mentions sample videos.

\section{Curriculum Development Context}
The New Testament is the main substantial text in the language; secular materials are scarce. A curriculum would need an agreed orthography and new texts, and the written-status disagreement should be checked with the community.
\end{document}